\noindent{\setfont{4}{1.5}\selectfont\bfseries {一、题目： \ProjectName}} \\[14bp]
\noindent{\setfont{4}{1.5}\selectfont\bfseries {二、指导老师对文献综述、开题报告、外文翻译的具体要求}}
\vspace{7bp}
\par
% ------ 不要动上面的 ------

文献综述不少于 3000 字，中外文献查阅篇数不少于 8 篇，其中外文文献建议为 3--5 篇。
开题报告不少于 3500 字。
外文翻译译文不少于 3000 字，译文须标注文章标题和原作者姓名，并要求提供外文原文及其出处。
\paragraph*{文献综述要求}
\begin{enumerate}
\item 简述与有限元方法、POD 方法以及缩减基方法（Reduced Basis Method）相关的教材和论文内容，掌握当前该领域的研究基础，并分析毕业论文研究内容与这些文献之间的联系与差异。
\item 分析 POD 方法以及本文拟研究的增量 POD 方法，简要说明这些方法的优缺点，并结合已有文献提出个人思考。
\item 概述所查阅参考文献的主要创新点，并指出其中存在的问题或尚未解决的问题。
\item 从所阅读的外文文献中选择一篇进行完整翻译，要求结构完整、语句通顺，并提供外文原文及其出处。
\end{enumerate}
\paragraph*{开题报告要求}
\begin{enumerate}
\item 简述模型降阶方法以及 POD 方法的应用背景。
\item 在文献综述基础上分析 POD 方法的缺陷，并说明研究增量 POD 方法的必要性。
\item 从理论分析与数值实验两个方面说明选题的可行性。
\item 阐述论文的主要研究内容。
\item 根据研究内容制定具体的实施计划与研究步骤。
\item 明确论文最终预期达到的研究成果。
\end{enumerate}

% ------ 不要动下面的（只改日期即可） ------
\newcommand{\foottext}[1]{{\setfont{-4}{1}\selectfont\bfseries {#1}}}
\vfill
\hfill \makebox[\widthof{\foottext{指导教师（签名）}}][l]{\foottext{指导教师（签名）}}\underline{\makebox[3cm][c]{\foottext{ }}}\par
\vspace{6bp}
\hfill \foottext{2026年3月6日}
\let\foottext\undefined